\documentclass[11pt,letterpaper]{article}

% ---------- Geometría y tipografía (igual que la guía) ----------
\usepackage[letterpaper,left=0.75in,right=0.75in,top=1.0in,bottom=0.75in,headheight=0pt,headsep=0pt,footskip=0pt]{geometry}
\usepackage{fontspec}
\setmainfont{Carlito}[
  Path=/nix/store/m7q6p26a3n2pyhnqhk0sbknfif2pbc8s-carlito-20130920/share/fonts/truetype/,
  Extension=.ttf,
  UprightFont=Carlito-Regular,
  BoldFont=Carlito-Bold,
  ItalicFont=Carlito-Italic,
  BoldItalicFont=Carlito-BoldItalic
]
\usepackage[spanish,es-noshorthands]{babel}
\usepackage{amsmath,amssymb}
\usepackage{xcolor}
\usepackage[most]{tcolorbox}
\usepackage{enumitem}
\usepackage{tabularx}
\usepackage{array}
\usepackage{titlesec}
\usepackage{needspace}
\usepackage{microtype}
\usepackage{graphicx}

% ---------- Paleta de la guía ----------
\definecolor{ink}{HTML}{2F3137}
\definecolor{muted}{HTML}{61666C}
\definecolor{accent}{HTML}{E11B22}
\definecolor{accentdark}{HTML}{A4121A}
\definecolor{boxgray}{HTML}{F3F5F7}
\definecolor{boxpink}{HTML}{FDEBEB}
\definecolor{boxcream}{HTML}{FFF4D6}
\definecolor{boxgreen}{HTML}{EAF4EF}
\definecolor{linegray}{HTML}{D9DDE1}

\color{ink}
\setlength{\parindent}{0pt}
\setlength{\parskip}{4pt}
\linespread{1.02}

% ---------- Encabezado y pie en todas las páginas ----------
\pagestyle{empty}
\AddToHook{shipout/background}{%
  \put(54,-64){\makebox(0,0)[bl]{\includegraphics[height=54pt]{LOGO-UA-color.jpg}}}%
  \put(558,-64){\makebox(0,0)[br]{\begin{minipage}[b]{250pt}\raggedleft
    {\color{muted}\fontsize{10}{12}\selectfont
     Cálculo II · Unidad 3\\
     Segundo semestre 2026\\
     Evaluación N.º 3}\end{minipage}}}%
  \put(54,-753.5){\color{accent}\rule{504pt}{0.5pt}}%
  \put(54,-757){\makebox(0,0)[tl]{\color{muted}\fontsize{10}{12}\selectfont Universidad Autónoma de Chile}}%
  \put(558,-757){\makebox(0,0)[tr]{\color{muted}\fontsize{10}{12}\selectfont\thepage}}%
}

% ---------- Cajas ----------
\newtcolorbox{cajagris}[1]{enhanced,frame hidden,arc=2mm,colback=boxgray,
  left=4mm,right=4mm,top=2.3mm,bottom=2.3mm,boxsep=0pt,coltext=ink,
  before skip=8pt,after skip=6pt,
  before upper={\textbf{#1}\par\vspace{3.5pt}}}
\newtcolorbox{cajapink}[1]{enhanced,frame hidden,arc=2mm,colback=boxpink,
  left=4mm,right=4mm,top=2.3mm,bottom=2.3mm,boxsep=0pt,coltext=ink,
  before skip=8pt,after skip=6pt,
  before upper={\textbf{#1}\par\vspace{3.5pt}}}
\newtcolorbox{cajacream}[1]{enhanced,frame hidden,arc=2mm,colback=boxcream,
  left=4mm,right=4mm,top=2.3mm,bottom=2.3mm,boxsep=0pt,coltext=ink,
  before skip=8pt,after skip=6pt,
  before upper={\textbf{#1}\par\vspace{3.5pt}}}
\newtcolorbox{cajagreen}[1]{enhanced,frame hidden,arc=2mm,colback=boxgreen,
  left=4mm,right=4mm,top=2.3mm,bottom=2.3mm,boxsep=0pt,coltext=ink,
  before skip=8pt,after skip=6pt,
  before upper={\textbf{#1}\par\vspace{3.5pt}}}

% ---------- Listas ----------
\newlist{squarelist}{itemize}{1}
\setlist[squarelist]{label={\color{ink}\scriptsize$\blacksquare$},
  leftmargin=13pt,labelsep=5pt,itemsep=2pt,topsep=2pt,parsep=0pt}

\newlist{checklist}{itemize}{1}
\setlist[checklist]{label={\color{ink}\small$\square$},
  leftmargin=13pt,labelsep=5pt,itemsep=2pt,topsep=2pt,parsep=0pt}

\newlist{preguntas}{enumerate}{1}
\setlist[preguntas]{label=\textbf{\color{accent}\arabic*.},
  leftmargin=22pt,labelsep=7pt,itemsep=12pt,topsep=4pt,parsep=2pt}

\newlist{partes}{enumerate}{1}
\setlist[partes]{label=\alph*),leftmargin=17pt,labelsep=5pt,
  itemsep=2pt,topsep=2pt,parsep=0pt}

% ---------- Títulos ----------
\titleformat{\section}{\color{accent}\fontsize{14.5}{17}\selectfont\bfseries}{\thesection.}{0.45em}{}
\titlespacing*{\section}{0pt}{14pt plus 2pt minus 2pt}{7pt}

\newcommand{\opciones}[5]{%
  \par\vspace{2pt}%
  \noindent
  \begin{tabularx}{\linewidth}{@{}>{\raggedright\arraybackslash}X@{\hspace{14pt}}>{\raggedright\arraybackslash}X@{}}
    \textbf{a)}~#1 & \textbf{b)}~#2\\[3pt]
    \textbf{c)}~#3 & \textbf{d)}~#4\\[3pt]
    \textbf{e)}~#5 & \null
  \end{tabularx}\par\vspace{-2pt}}

\newcommand{\ejercicio}[2]{%
  \par\needspace{6\baselineskip}%
  \noindent{\bfseries\color{ink}Ejercicio #1: #2}\par}

\newcommand{\rotulo}[1]{\par\vspace{4pt}{\color{ink}\bfseries\fontsize{10.5}{12}\selectfont #1}\par\vspace{1pt}}

\newcount\nlineas
\newcommand{\lineas}[1]{%
  \par\vspace{1pt}%
  \nlineas=#1\relax
  \loop\ifnum\nlineas>0
    \noindent{\color{linegray}\rule{\linewidth}{0.35pt}}\par\vspace{19pt}%
    \advance\nlineas by -1
  \repeat
  \vspace{-10pt}}

\begin{document}

% ================= PORTADA / BLOQUE INICIAL =================
{\color{ink}\bfseries\fontsize{17.3}{20}\selectfont Taller de refuerzo}\par\vspace{1pt}
{\color{ink}\fontsize{14.7}{18}\selectfont Sucesiones y series}\par\vspace{9pt}

\noindent
\mbox{\textbf{Nombre:}\ \rule[0.5ex]{0.54\linewidth}{0.4pt}}%
\hfill
\mbox{\textbf{Fecha:}\ \rule[0.5ex]{0.22\linewidth}{0.4pt}}\par\vspace{2pt}

% ================= SECCIÓN 1 =================
\section{Preguntas de alternativa conceptuales}

Responde sin calculadora y encierra en un círculo una sola alternativa por pregunta. Las respuestas incorrectas no descuentan.

\begin{preguntas}
\needspace{7\baselineskip}
\item La sucesión $a_n=3+\dfrac{2}{n}$, con $n\ge 1$:
  \opciones{converge a $2$ y es creciente.}
           {converge a $3$ y es decreciente.}
           {diverge, porque $\frac{1}{n}$ no está definida en $n=0$.}
           {converge a $3$ y es creciente.}
           {converge a $0$, porque $\frac{2}{n}\to 0$.}

\needspace{7\baselineskip}
\item Sea una sucesión geométrica con $a_1=5$ y razón $r=-\frac{1}{3}$. ¿Cuál es el valor de $a_4$?
  \opciones{$\dfrac{5}{27}$}
           {$-\dfrac{5}{27}$}
           {$-\dfrac{5}{81}$}
           {$\dfrac{5}{81}$}
           {$-45$}

\needspace{7\baselineskip}
\item Una serie geométrica con $a_1\ne 0$ converge si y solo si:
  \opciones{$r<1$.}
           {$|r|<1$.}
           {$|r|\le 1$.}
           {$r=0$.}
           {que $r$ es positivo.}

\needspace{6\baselineskip}
\item El límite de la sucesión $a_n=\dfrac{2n^2+n}{n^2+3}$ es:
  \opciones{$0$}
           {$1$}
           {$2$}
           {$3$}
           {$+\infty$}

\newpage
\needspace{9\baselineskip}
\item Para aplicar el criterio integral a $\sum a_n$ se busca $f$ tal que $a_n=f(n)$. ¿Qué condiciones son suficientes en $[N,\infty)$?
  \opciones{$f$ continua y positiva.}
           {$f$ continua, positiva y decreciente.}
           {$f$ derivable y acotada.}
           {$f$ positiva y con límite $0$ cuando $x\to\infty$.}
           {$f$ creciente y no negativa.}

\needspace{8\baselineskip}
\item Si $\lim_{n\to\infty} a_n=0$, entonces:
  \opciones{la serie $\sum a_n$ converge.}
           {la serie $\sum a_n$ diverge.}
           {la serie converge solo si es geométrica.}
           {no se puede concluir nada sobre la convergencia de $\sum a_n$.}
           {la suma de la serie es $0$.}

\needspace{7\baselineskip}
\item Al aplicar el criterio del cociente a una serie de términos no nulos se obtiene $L=1$. ¿Qué se concluye?
  \opciones{La serie converge.}
           {La serie diverge.}
           {La serie converge absolutamente.}
           {El criterio no decide; se debe usar otro criterio.}
           {La suma de la serie es $1$.}

\needspace{8\baselineskip}
\item Para la serie $\displaystyle\sum_{n=1}^{\infty}\frac{4^n}{n!}$, el criterio del cociente entrega:
  \opciones{$L=0<1$; converge.}
           {$L=4>1$; diverge.}
           {$L=\frac{1}{4}<1$; converge.}
           {$L=1$; no se decide.}
           {$L=+\infty$; diverge.}

\needspace{6\baselineskip}
\item Si la suma parcial de una serie es $S_N=3-\dfrac{1}{N}$, entonces la suma de la serie es:
  \opciones{$0$}
           {$1$}
           {$2$}
           {$3$}
           {divergente.}
\end{preguntas}

% ================= SECCIÓN 2 =================
\clearpage
\section{Resuelve las siguientes sucesiones y series}

En cada ejercicio explica por escrito cada paso: qué criterio usas, por qué se cumplen sus hipótesis y qué concluyes. No basta con indicar el resultado. Salvo que se indique otra cosa, las sucesiones comienzan en $n=1$.

\ejercicio{1}{Sucesión geométrica contextualizada}\par\vspace{4pt}
Un equipo transmite $800$ MB en la primera etapa. En cada etapa posterior transmite el $30$\% de la información enviada en la etapa anterior. El proceso continúa indefinidamente.
\begin{partes}
\item Escribe los cuatro primeros términos y el término general $a_n$.
\item Determina si la sucesión es monótona y acotada. Calcula $\displaystyle\lim_{n\to\infty}a_n$ e interpreta el resultado en el contexto.
\item Expresa la información total transmitida como una serie geométrica y calcula su suma. Indica la unidad.
\end{partes}

\rotulo{Desarrollo}
\lineas{12}

\clearpage
\ejercicio{2}{Límite de una sucesión con raíz}\par\vspace{4pt}
Considera la sucesión
\[
a_n=\sqrt{n^2+4n}-n.
\]
\begin{partes}
\item Multiplica y divide por el conjugado $\sqrt{n^2+4n}+n$.
\item Simplifica la expresión y calcula $\displaystyle\lim_{n\to\infty}a_n$, justificando cada paso.
\item Decide si la sucesión es convergente. Si lo es, explica qué significa el valor obtenido.
\end{partes}

\rotulo{Desarrollo}
\lineas{12}

\clearpage
\ejercicio{3}{Criterio integral}\par\vspace{4pt}
Estudia la convergencia de la serie
\[
\sum_{n=3}^{\infty}\frac{n}{n^2+9}.
\]
\begin{partes}
\item Define $f(x)=\dfrac{x}{x^2+9}$ y verifica continuidad, positividad y decrecimiento en $[3,\infty)$. Usa $f'(x)$ para justificar el decrecimiento.
\item Escribe la integral impropia como un límite y calcúlala con la sustitución $u=x^2+9$.
\item Concluye sobre la convergencia o divergencia de la serie. Explica por qué $a_n\to 0$ no es suficiente para afirmar que converge.
\end{partes}

\rotulo{Desarrollo}
\lineas{12}

% ================= SECCIÓN 3 =================
\clearpage
\section{Aplicación física: una sucesión que se detiene}

Resuelve la situación y justifica cada paso por escrito.

\begin{cajagris}{Situación}
Una pelota se deja caer desde $2$ m. Después del primer rebote alcanza el $60$\% de la altura inicial; después de cada rebote siguiente alcanza el $60$\% de la altura del rebote anterior. Sea $h_n$ la altura máxima, en metros, alcanzada después del $n$-ésimo rebote.
\end{cajagris}

\begin{partes}
\item Calcula $h_1$, $h_2$ y $h_3$. Escribe el término general $h_n$.
\item Determina si $(h_n)$ es monótona y acotada. Calcula $\displaystyle\lim_{n\to\infty}h_n$.
\item Concluye si la sucesión es convergente e interpreta físicamente el límite.
\end{partes}

\rotulo{Desarrollo}
\lineas{9}
\rotulo{Interpretación física:}
\lineas{3}

\end{document}
